(1) 取斜面向下为正、接触处静摩擦力 \(f\) 的代数分量也取向下。纯滚动的质心加速度 \(a=\ddot x\) 与角加速度满足 \(|\alpha_{\rm rot}|=|a|/R\)。取与向下滚动相容的转向为正，转动方程为 \(I a/R=-fR\)，因此 \(f=-I a/R^2=-\lambda ma\)。质心平动给
\[
m\ddot x=mg\sin\theta-kx+f
\]
所以
\[
m(1+\lambda)\ddot x+kx=mg\sin\theta .
\]
摩擦力随加速度变号：向下加速时指向斜面上方，回程向上加速时则指向斜面下方。

(2) 静平衡 \(x_{\rm eq}=mg\sin\theta/k\)，令 \(\xi=x-x_{\rm eq}\)，有 \(\ddot\xi+\omega^2\xi=0\)，其中
\[
\omega=\sqrt{\frac{k}{m(1+\lambda)}}.
\]
初速度零，因此
\[
x(t)=x_{\rm eq}+(x_0-x_{\rm eq})\cos\omega t,\qquad
\dot x(t)=-\omega(x_0-x_{\rm eq})\sin\omega t .
\]
纯滚动的转动惯量降低频率，但不改变静平衡位置。

(3) 平动与转动动能分别是 \(K_{\rm trans}=m\dot x^2/2\) 与 \(K_{\rm rot}=I(\dot x/R)^2/2=\lambda m\dot x^2/2\)；势能包括 \(U_s=kx^2/2\) 与重力势能 \(U_g=-mgx\sin\theta\)。整理平方，得到
\[
E=\frac{m(1+\lambda)}2\dot x^2
+\frac{k}{2}(x-x_{\rm eq})^2
-\frac k2x_{\rm eq}^2 .
\]
其时间导数为 \(\dot x[m(1+\lambda)\ddot x+k(x-x_{\rm eq})]=0\)，所以静摩擦在纯滚动接触点不做功、总能量恒定。

(4) 由 \(f=-\lambda m\ddot x\) 与运动解，\(|f|_{\max}=\lambda m\omega^2|x_0-x_{\rm eq}|=\lambda k|x_0-x_{\rm eq}|/(1+\lambda)\)。法向支持力 \(N=mg\cos\theta\)，因此整段无滑动的充分且在库仑静摩擦模型中必要的条件是
\[
\mu_s mg\cos\theta\ge\frac{\lambda k}{1+\lambda}|x_0-x_{\rm eq}|.
\]
当 \(\lambda\to0^+\)，所需静摩擦趋零、\(\omega\to\sqrt{k/m}\)；当 \(\theta\to0\)，\(x_{\rm eq}\to0\)，退化为水平纯滚动弹簧振子。
